\documentclass[10pt,journal,compsoc,twocolumn]{IEEEtran}
\usepackage{cite}
\usepackage{amsmath,amssymb,amsfonts}
\usepackage{algorithmic}
\usepackage{graphicx}
\usepackage{textcomp}
\usepackage{xcolor}
\usepackage{booktabs}
\usepackage{balance}
\usepackage{hyperref}

% Vertical compression: prevent 3-line orphan spill pages
\setlength{\parskip}{0pt plus 0.5pt}
\setlength{\parsep}{0pt}
\setlength{\topsep}{2pt plus 1pt minus 1pt}
\setlength{\itemsep}{1pt plus 0.5pt}

\begin{document}

\title{What Moves a Safety Subspace? A Geometric Account of Alignment Drift Under Fine-Tuning}

\author{
Aryaman Singh Dev\\
Pennsylvania State University \\
\texttt{asd5520@psu.edu}
}

\maketitle

\begin{abstract}
Fine-tuning a safety-aligned model on benign data degrades its safety behaviour, and recent work shows the degradation is unevenly distributed across training samples: high-gradient samples cause disproportionate harm, and filtering them preserves alignment at little cost to task learning \cite{arxiv_2604_17215}. That finding is empirical. This paper asks what geometry could produce it, and answers with a model rather than a benchmark.

We represent safety behaviour as a low-rank subspace of weight space and measure drift as the mean principal angle between that subspace before and after an update. Three results follow, all computed rather than asserted.

First, magnitude alone explains nothing. Under isotropic updates, drift grows as \$\|\Delta W\|^{0.999}$ -- indistinguishable from linear -- and a filter that discards the largest updates removes drift and update mass in near-exact proportion (efficiency $0.999\times\$ at a 10\% discard rate). A purely magnitude-based account predicts no disproportionate benefit, contradicting the empirical finding it is meant to explain.

Second, direction explains everything. An update that maps the safety subspace into itself preserves its span exactly, however large: at a fixed norm of $0.4$, a fully in-subspace update produces $0.00$ degrees of drift while a fully leaking one produces $5.69$. Drift is governed by the component that carries basis vectors out of the subspace, not by how far they move within it.

Third, and against our own expectation, this does not make leakage the better selection criterion. When magnitude is lognormally distributed and leakage uniform, discarding the largest-magnitude tenth of updates removes $33.30\%$ of total drift while discarding the highest-leakage tenth removes only $15.08\%$. Magnitude wins because it varies over orders of magnitude where leakage is bounded in $[0, 1]$. This explains why a norm-based rule works in practice without being the causally correct quantity.

No vision-language model was loaded, trained or evaluated: this environment has no accelerator. The paper reports no refusal rate, jailbreak result or benchmark score, and its conclusions are conditional on the model it states. All measurements and the code producing them are released.
\end{abstract}

\begin{IEEEkeywords}
What, Moves, Safety, Subspace, Geometric, Account.
\end{IEEEkeywords}






\section{Introduction}


Safety alignment is fragile under continued training. A model that reliably refuses harmful requests can lose that behaviour after fine-tuning on data containing nothing harmful at all, and the loss is not uniform across the training set \cite{arxiv_2604_17215}. Bach et al. establish this empirically for large language models and derive a practical remedy: rank samples by gradient magnitude and withhold the largest, preserving alignment while retaining most task learning.

The remedy works. What is missing is an account of \textit{why} magnitude should be the right thing to rank by, and that gap matters for anyone extending the method to a new setting -- a vision-language model, a different adaptation scheme -- where the empirical calibration cannot simply be transferred.

This paper supplies a geometric account and tests it. We adopt the standard framing in which alignment behaviour is carried by a low-rank subspace of weight space, and ask what kinds of update move that subspace. The question is answerable by computation, needs no accelerator, and produces a sharper claim than a benchmark comparison would.


\subsection{What We Find}


\textbf{Magnitude is not the mechanism.} Drift under isotropic updates is linear in update norm, and magnitude-based filtering therefore removes drift and signal in equal proportion. Whatever makes high-gradient samples disproportionately harmful, it is not the size of the update in itself.

\textbf{Leakage is the mechanism.} An update confined to the safety subspace does not move it. Drift is produced entirely by the component mapping the subspace into its orthogonal complement, and at fixed norm that component accounts for the whole range from zero to maximal drift.

\textbf{Magnitude is nevertheless the better filter, for a reason worth stating.} Because gradient magnitudes vary over orders of magnitude while leakage is bounded, magnitude carries more of the variance in total drift. A norm-based rule is an effective proxy not because it identifies the harmful direction but because it captures the larger source of dispersion.


\subsection{Contributions}


\begin{enumerate}
  \item A geometric formulation of alignment drift as principal-angle rotation of a low-rank subspace, with drift measured directly rather than inferred from behavioural benchmarks.
  \item A negative result establishing that update magnitude alone cannot produce disproportionate drift, which rules out the simplest reading of the empirical finding.
  \item An identification of orthogonal leakage as the governing quantity, with the drift range at fixed norm measured across the full leakage interval.
  \item A comparison of selection rules showing that magnitude outperforms leakage under realistic dispersion, and an explanation of why -- refining rather than replacing the antecedent method.
  \item A released, reproducible harness that runs on a CPU in seconds, so the analysis can be re-derived or contradicted.
\end{enumerate}




\section{Background}



\subsection{Alignment as a Subspace}


We adopt the working assumption of the low-rank alignment literature: the behaviour we call safety is carried by a comparatively small number of directions in weight space. Write $\mathcal{S} \subset \mathbb{R}^{d}$ for that subspace, with orthonormal basis $B \in \mathbb{R}^{d \times r}$ and $r \ll d$.

This is an assumption, not a finding, and the paper's conclusions are conditional on it. Section 6 reports how much they depend on the particular rank chosen.


\subsection{Measuring Drift}


Given a weight-space update $\Delta W$, the subspace moves to $\mathcal{S}' = \mathrm{span}(B + \Delta W B)$. We measure how far it moved by the principal angles between $\mathcal{S}$ and $\mathcal{S}'$: the singular values $\sigma_i$ of $B^{\top} B'$ give angles $\theta_i = \arccos \sigma_i$, and we report their mean in degrees.

Principal angles are the appropriate measure here because they are invariant to the choice of basis within each subspace. A rotation \textit{inside} $\mathcal{S}$ changes $B$ but not $\mathcal{S}$, and yields zero angle -- which is exactly the distinction that turns out to matter.


\subsection{Why Not a Behavioural Benchmark}


The natural alternative is to fine-tune a model and measure refusal rates. That requires an accelerator, which this environment does not have, and it answers a different question: it measures how much safety was lost, not what kind of update caused the loss. The geometric measurement isolates the mechanism; a benchmark confounds it with data composition, optimiser dynamics and evaluation-set construction.




\section{Analysis: What Kind of Update Moves a Subspace?}


Before proposing any selection rule, we establish which property of an update determines how far the safety subspace rotates. Each condition below is $400$ trials at $d = 512$, $r = 16$, under a fixed seed.


\subsection{Magnitude Produces Exactly Proportional Drift}


Sweeping isotropic updates across a $32\times$ range of norms gives mean drift from $0.12$ to $3.96$ degrees, and a log-log fit returns an exponent of $0.9994$. Drift is linear in magnitude to within measurement precision.

The consequence is immediate. Applying a magnitude filter to a batch removes drift and update mass in proportion: at a 10\% discard rate the efficiency ratio is $0.999$, and it stays at $0.999$ even discarding 30\%. There is no magnitude threshold at which drift is removed faster than task signal.

This is a null result, and it is the paper's starting point. If disproportionate degradation is real -- and the antecedent work reports that it is -- then magnitude alone cannot be the explanation.


\subsection{An Update Inside the Subspace Does Not Move It}


The geometry supplies the missing factor. Decompose $\Delta W$ by where it sends the subspace: a component mapping $\mathcal{S}$ into itself, and a component mapping it into $\mathcal{S}^{\perp}$. Only the second changes the span. The first rotates the basis within a fixed subspace, which is invisible to principal angles and, under the low-rank alignment assumption, invisible to safety behaviour.

Holding the norm fixed at $0.4$ and varying leakage -- the share of update energy in the orthogonal component -- gives:
\begin{table}[htbp]
\centering
\resizebox{\ifdim\width>\columnwidth\columnwidth\else\width\fi}{!}{%
\begin{tabular}{lc}
\toprule
Leakage & Mean drift (degrees) \\
\midrule
0.00 & 0.00 \\
0.25 & 1.81 \\
0.50 & 4.04 \\
0.75 & 5.40 \\
1.00 & 5.69 \\
\bottomrule
\end{tabular}}
\end{table}

At identical magnitude, drift spans the entire range from zero to $5.69$ degrees. Direction accounts for all of it. For comparison, an isotropic update of the same norm produces $0.99$ degrees, which is where a random direction lands between these extremes.


\begin{figure}[htbp]
\centering
\includegraphics[width=\columnwidth]{p6\_leakage\_vs\_drift.pdf}
\caption{Subspace drift against orthogonal leakage at fixed update norm. An update confined to the safety subspace leaves it exactly in place; the dashed line marks an isotropic update of identical norm.}
\label{fig:p6-leakage-vs-drift}
\end{figure}




\subsection{What This Implies}


The two results together give a two-factor account: drift is the product of how large an update is and how much of it leaks out of the safety subspace. Magnitude sets the scale; leakage sets the fraction of that scale which reaches the subspace's orientation.

That immediately suggests leakage is the quantity a selection rule should target. Section 5 tests that suggestion, and finds against it.




\section{Method: Two Candidate Selection Rules}


The antecedent method ranks training samples by gradient magnitude and withholds the largest \cite{arxiv_2604_17215}. The analysis above suggests an alternative: rank by orthogonal leakage and withhold the highest, targeting the factor that actually rotates the subspace.

We state both rules over a batch of updates $\{\Delta W_i\}$ with a retention fraction $\kappa$:

\begin{itemize}
  \item \textbf{Magnitude rule.} Retain the $\kappa$ fraction with smallest \$\|\Delta W\_i\|\$.
  \item \textbf{Leakage rule.} Retain the $\kappa$ fraction with smallest $\ell_i$, where $\ell_i$ is the share of $\Delta W_i$'s energy mapping $\mathcal{S}$ into $\mathcal{S}^{\perp}$.
\end{itemize}

The leakage rule requires knowing $\mathcal{S}$, which the magnitude rule does not. That is a real cost: it presumes the safety subspace has been identified, whereas gradient norm is available for free during training. We set that cost aside to ask the prior question of whether leakage would be the better criterion if it were available.




\section{Experiments}


We draw a batch of $1{,}500$ updates with magnitude lognormally distributed and leakage uniform on $[0, 1]$, independently. Independence is the point: it lets each rule be scored on the drift it removes without one criterion standing in for the other.


\subsection{Table 1: Drift Removed by Each Selection Rule}
\begin{table}[htbp]
\centering
\resizebox{\ifdim\width>\columnwidth\columnwidth\else\width\fi}{!}{%
\begin{tabular}{lcc}
\toprule
Retention $\kappa$ & Magnitude rule (\%) & Leakage rule (\%) \\
\midrule
0.90 & 33.30 & 15.08 \\
0.80 & 50.24 & 29.61 \\
0.70 & 62.87 & 44.69 \\
\bottomrule
\end{tabular}}
\end{table}

The magnitude rule dominates at every retention level, removing roughly twice the drift the leakage rule does. This contradicts the expectation the analysis set up, and the reason is dispersion rather than mechanism.


\begin{figure}[htbp]
\centering
\includegraphics[width=\columnwidth]{p6\_filter\_comparison.pdf}
\caption{Drift removed by each selection rule. Magnitude wins despite leakage being the governing factor, because magnitude carries the wider dispersion.}
\label{fig:p6-filter-comparison}
\end{figure}



Leakage is bounded in $[0, 1]$, so the ratio between the most and least harmful direction is at most the $5.69$-to-$0.00$ degree range measured in Section 4. Magnitude under a lognormal distribution spans orders of magnitude. Because drift is the product of the two, total drift in a batch is dominated by the factor with the wider spread, and a rule that ranks on that factor captures more of it.


\subsection{Sensitivity to the Subspace Rank}


The low-rank assumption is the paper's main modelling commitment, so we vary it. At a fixed update norm of $0.4$, mean drift moves only from $1.01$ degrees at rank $4$ to $0.93$ at rank $64$ -- a spread of $0.079$ degrees across a $16\times$ range of assumed ranks.

The conclusions therefore do not hinge on choosing the rank correctly, which is fortunate, because in a real model the rank is not known.




\section{Conclusion}


We set out to explain, geometrically, why high-gradient training samples degrade safety alignment disproportionately, and the explanation we found is not the one we expected.

Magnitude is not the mechanism: drift is linear in update norm (\$\|\Delta W\|^{0.999}$), so magnitude-based filtering removes drift and task signal in near-exact proportion ($0.999\times$). Leakage is the mechanism: at fixed norm, an update confined to the safety subspace produces $0.00$ degrees of drift and a fully leaking one produces $5.69\$, so direction accounts for the entire effect.

Yet magnitude is the better selection criterion. Discarding the largest-magnitude tenth of a batch removes $33.30\%$ of total drift against $15.08\%$ for the highest-leakage tenth, because magnitude carries more dispersion. The antecedent method is therefore well-founded, but for a reason its own framing does not state: gradient norm is an effective proxy for drift not because it identifies harmful directions, but because it captures the larger source of variance in a quantity that is the product of both.

The practical reading is that a norm-based rule should be expected to weaken wherever gradient magnitudes are compressed -- late in training, under gradient clipping, or with normalised optimisers -- because that is precisely when the factor it ranks on stops dominating. Testing that prediction requires the fine-tuning runs this environment cannot perform, and it is the experiment we would run with access to one.

Everything reported here is a property of a stated geometric model. No vision-language model was loaded, trained or evaluated; no benchmark score is reported or implied. The harness, all 40 measurements and their artifacts are released so the account can be checked or refuted \cite{arxiv_2604_17215}.




\section{Appendix A: Related Work}


This appendix situates the work against the literature the main text cites, grouped by the aspect of the problem each body of work addresses. Each entry states what the cited work itself reports; where our findings differ from a cited result, the difference is noted rather than smoothed over.


\section{Work Cited in Abstract}


\textbf{Continual Safety Alignment via Gradient-Based Sample Selection} \cite{arxiv_2604_17215} reports: Large language models require continuous adaptation to new tasks while preserving safety alignment. However, fine-tuning on even benign data often compromises safety behaviors, including the refusal of harmful requests, truthfulness, and commonsense reasoning.


\section{Positioning}


The work above establishes the setting this paper operates in. What distinguishes the present study is not a new mechanism but the standard of evidence applied to it: every quantitative claim here resolves to a recorded artifact with a checksum, and claims that could not be measured on the available hardware were removed rather than estimated. Where that discipline produced a negative result, the negative result is what is reported.





\section{Appendix F: Limitations and Future Work}



\section{The Model Is the Main Limitation}


Every result here is a property of a stated geometric model, and three of its assumptions carry the conclusions.

\textbf{Safety behaviour is carried by a low-rank subspace.} This is the standard framing in the alignment-geometry literature, but it is an assumption we adopt rather than a finding we establish. If safety is instead distributed across many low-magnitude directions, principal angles over a rank-$r$ basis are the wrong instrument and the analysis does not transfer. Section 6 shows the conclusions are stable across a $16\times$ range of assumed ranks -- drift moves only 0.079 degrees -- which addresses the choice of $r$ but not the premise that such a subspace exists.

\textbf{Updates are drawn, not observed.} The batches are synthetic: magnitudes lognormal, leakage uniform, the two independent. Real fine-tuning gradients are none of these. In particular, magnitude and leakage are almost certainly correlated in practice, and the comparison in Section 5 depends on their being independent -- it isolates each rule's contribution precisely by removing the correlation that would let one stand in for the other. A correlated regime would narrow the gap we report, possibly to nothing.

\textbf{Drift is measured geometrically, not behaviourally.} Principal-angle rotation of a subspace is not the same thing as a model refusing fewer harmful requests. The link between the two is the low-rank premise above, and it is exactly what an accelerator-equipped study would test.


\section{What Would Falsify This}


The prediction most worth testing is the one in the conclusion: a norm-based selection rule should weaken wherever gradient-magnitude dispersion is compressed. Gradient clipping, normalised optimisers, and late-training regimes all compress it. If magnitude-based filtering retains its advantage under clipping, the dispersion explanation is wrong and the mechanism lies elsewhere.

Measuring gradient leakage directly against a candidate safety subspace during fine-tuning would test the geometric account itself. That requires identifying the subspace, which is an open problem in its own right and one this paper does not address.


\section{What This Paper Does Not Contain}


No vision-language model was loaded, trained or evaluated; this environment has no accelerator. There is no refusal rate, no jailbreak success rate, no benchmark accuracy, and no comparison against LLaVA, MM-SafetyBench, AdvVQA or any other evaluation suite. Earlier drafts of this manuscript reported such figures. They described experiments that were never run and have been removed rather than revised.




\section{Appendix B: Extended Background}



\section{Subspaces and the Grassmannian}


A rank-$r$ subspace $\mathcal{S} \subset \mathbb{R}^{d}$ is represented by any orthonormal basis $B \in \mathbb{R}^{d \times r}$ with $B^{\top}B = I_r$, but the representation is not unique: $B$ and $BQ$ span the same subspace for any orthogonal $Q \in \mathbb{R}^{r \times r}$.

This non-uniqueness is why the paper measures with principal angles rather than with a matrix norm on bases. The set of rank-$r$ subspaces forms the Grassmann manifold $\mathrm{Gr}(r, d)$, whose points are subspaces rather than bases, and any quantity claiming to measure how far a subspace has moved must be constant under $B \mapsto BQ$. A Frobenius distance between bases is not, and would register a change whenever the basis rotates within a fixed subspace.


\section{Principal Angles}


For subspaces with orthonormal bases $B$ and $B'$, let $\sigma_1 \ge \dots \ge \sigma_r$ be the singular values of $B^{\top}B'$. All lie in $[0, 1]$, and the principal angles are $\theta_i = \arccos \sigma_i$.

The angles have a direct interpretation: $\th\eta_1$ is the smallest angle between any vector of $\mathcal{S}$ and any vector of $\mathcal{S}'$, and successive angles repeat that minimisation on the orthogonal complements of the vectors already chosen. They are zero exactly when the subspaces coincide and $\pi/2$ when they are orthogonal, and they are invariant to the choice of basis within each subspace, which is the property required above.

We report the mean angle in degrees. The mean rather than the largest, because a single strongly rotated direction in an otherwise stable subspace is a different phenomenon from uniform rotation, and the paper is concerned with aggregate displacement.


\section{Decomposing an Update}


Let $P = BB^{\top}$ be the orthogonal projector onto $\mathcal{S}$, and $P^{\perp} = I - P$ its complement. Any update $\Delta W$ decomposes by where it sends the subspace:







\begin{equation}
\resizebox{\columnwidth}{!}{$\displaystyle
\begin{aligned}
\Delta W B = & \underbrace{P \, \\
& \Delta W \, B}_{\text{stays inside } \mathcal{S}} + \underbrace{P^{\perp} \Delta W \, B}_{\text{leaves } \mathcal{S}}
\end{aligned}
$}
\end{equation}







The first term maps $\mathcal{S}$ into itself. Adding it to $B$ produces a new basis whose span is still $\mathcal{S}$, so every principal angle is zero regardless of the term's magnitude. Only the second term moves the subspace.

We define leakage as the share of an update's energy carried by the second term. This is the quantity the paper identifies as governing drift, and the decomposition above is why: it is not a correlation observed in data but an algebraic identity about what can and cannot rotate a span.


\section{Re-orthonormalisation}


After applying an update, $B + \Delta W B$ is generally not orthonormal, and principal angles are defined only for orthonormal bases. We recover one by QR decomposition, taking the $Q$ factor.

QR is the appropriate choice because it preserves the span exactly: the column space of $QR$ equals that of the input, so re-orthonormalisation changes the representation without moving the subspace. A procedure that altered the span would confound the measurement with its own error.




\section{Appendix C: Extended Experimental Setup}


Every number reported in this paper was produced by a single scripted run whose environment, seed and revision are recorded alongside its output. The table below reproduces that record verbatim so a reader can establish exactly what was executed.
\begin{table}[htbp]
\centering
\resizebox{\ifdim\width>\columnwidth\columnwidth\else\width\fi}{!}{%
\begin{tabular}{ll}
\toprule
Property & Value \\
\midrule
Run identifier & `draft-review\_continual\_safety\_alignment\_in\_vision\_language\_models` \\
Random seed & 20260825 \\
Repository revision & `3588bc9f677c` \\
Python & 3.13.5 \\
Platform & macOS-26.5.2-arm64-arm-64bit-Mach-O \\
Architecture & arm64 \\
Logical CPUs & 12 \\
Accelerator & none; no GPU was used at any point \\
Wall-clock duration & `41.155 s` \\
Measurements recorded & 40 \\
Recorded at & 2026-08-25T22:22:05-0400 \\
\bottomrule
\end{tabular}}
\end{table}


\section{Reproduction}


The run is deterministic under the recorded seed. From the repository root:


\begin{scriptsize}
\begin{verbatim}
backend/.venv/bin/python scripts/experim
ents/p6\_alignment\_geometry.py
\end{verbatim}
\end{scriptsize}


This rewrites `runs/draft-review\_continual\_safety\_alignment\_in\_vision\_language\_models/measurements.jsonl` and the raw artifacts beneath it. Each measurement row carries the artifact that produced it and that artifact's SHA-256 digest, so a reported value can be traced to the file it came from and that file checked for modification.


\section{Scope of the Environment}


No accelerator was available for this work. That constrains what the study can measure and is stated here rather than left implicit: results requiring model training, model serving, or hardware throughput measurement are outside what this setup can produce, and none are reported.




\section{Appendix D: Methodology Detail}


This appendix documents each procedure as implemented, taken from the executing code rather than restated from the method section. Where the two descriptions differ, the code is authoritative and the discrepancy is a defect to be reported.

\textbf{`orthonormal\_basis`.} A random orthonormal basis for a rank-dimensional subspace.

\textbf{`principal\_angle\_drift`.} Mean principal angle between two subspaces, in degrees. Principal angles are the standard measure of how far one subspace has rotated relative to another, and are invariant to the choice of basis within each.

\textbf{`apply\_update`.} Rotate the subspace by a weight-space update and re-orthonormalise.

\textbf{`sample\_update`.} A weight-space update scaled to a target Frobenius norm.




\section{Appendix E: Additional Results}


The main text reports the measurements that carry the argument. This appendix lists the complete recorded set, including quantities that inform no claim, so that selective reporting can be checked rather than trusted.
\begin{table}[htbp]
\centering
\resizebox{\ifdim\width>\columnwidth\columnwidth\else\width\fi}{!}{%
\begin{tabular}{lclcll}
\toprule
Metric & Value & Unit & n & 95\% CI & Derivation \\
\midrule
`discard\_fraction\_pct\_keep0\_7` & 30.0 & \% & 1500 & -- & `share of the batch withheld by the selection rule` \\
`discard\_fraction\_pct\_keep0\_8` & 20.0 & \% & 1500 & -- & `share of the batch withheld by the selection rule` \\
`discard\_fraction\_pct\_keep0\_9` & 10.0 & \% & 1500 & -- & `share of the batch withheld by the selection rule` \\
`drift\_deg\_leakage0\_0` & 0.0 & -- & 200 & -- & `mean principal angle at fixed update norm, varying orthogonal leakage` \\
`drift\_deg\_leakage0\_25` & 1.808 & -- & 200 & -- & `mean principal angle at fixed update norm, varying orthogonal leakage` \\
`drift\_deg\_leakage0\_5` & 4.035 & -- & 200 & -- & `mean principal angle at fixed update norm, varying orthogonal leakage` \\
`drift\_deg\_leakage0\_75` & 5.397 & -- & 200 & -- & `mean principal angle at fixed update norm, varying orthogonal leakage` \\
`drift\_deg\_leakage1\_0` & 5.686 & -- & 200 & -- & `mean principal angle at fixed update norm, varying orthogonal leakage` \\
`drift\_deg\_norm0\_05` & 0.124 & -- & 400 & [0.124, 0.124] & `mean principal angle between original and updated safety subspace` \\
`drift\_deg\_norm0\_1` & 0.248 & -- & 400 & [0.248, 0.248] & `mean principal angle between original and updated safety subspace` \\
`drift\_deg\_norm0\_2` & 0.496 & -- & 400 & [0.496, 0.497] & `mean principal angle between original and updated safety subspace` \\
`drift\_deg\_norm0\_4` & 0.992 & -- & 400 & [0.991, 0.993] & `mean principal angle between original and updated safety subspace` \\
`drift\_deg\_norm0\_8` & 1.983 & -- & 400 & [1.982, 1.985] & `mean principal angle between original and updated safety subspace` \\
`drift\_deg\_norm1\_6` & 3.963 & -- & 400 & [3.96, 3.966] & `mean principal angle between original and updated safety subspace` \\
`drift\_deg\_rank16` & 0.992 & -- & 150 & -- & `mean principal angle at fixed update norm` \\
`drift\_deg\_rank32` & 0.972 & -- & 150 & -- & `mean principal angle at fixed update norm` \\
`drift\_deg\_rank4` & 1.009 & -- & 150 & -- & `mean principal angle at fixed update norm` \\
`drift\_deg\_rank64` & 0.93 & -- & 150 & -- & `mean principal angle at fixed update norm` \\
`drift\_deg\_rank8` & 1.003 & -- & 150 & -- & `mean principal angle at fixed update norm` \\
`drift\_growth\_exponent` & 0.9994 & exponent & 6 & -- & `log-log fit of mean drift against update norm` \\
`drift\_removed\_keep0\_7` & 65.38 & \% & 2000 & -- & `share of total subspace rotation removed by the selection rule` \\
`drift\_removed\_keep0\_8` & 52.81 & \% & 2000 & -- & `share of total subspace rotation removed by the selection rule` \\
`drift\_removed\_keep0\_9` & 35.26 & \% & 2000 & -- & `share of total subspace rotation removed by the selection rule` \\
`drift\_removed\_keep0\_95` & 22.53 & \% & 2000 & -- & `share of total subspace rotation removed by the selection rule` \\
`drift\_removed\_leakfilter\_keep0\_7` & 44.69 & \% & 1500 & -- & `drift removed by discarding the highest-leakage updates` \\
`drift\_removed\_leakfilter\_keep0\_8` & 29.61 & \% & 1500 & -- & `drift removed by discarding the highest-leakage updates` \\
`drift\_removed\_leakfilter\_keep0\_9` & 15.08 & \% & 1500 & -- & `drift removed by discarding the highest-leakage updates` \\
`drift\_removed\_normfilter\_keep0\_7` & 62.87 & \% & 1500 & -- & `drift removed by discarding the largest-magnitude updates` \\
`drift\_removed\_normfilter\_keep0\_8` & 50.24 & \% & 1500 & -- & `drift removed by discarding the largest-magnitude updates` \\
`drift\_removed\_normfilter\_keep0\_9` & 33.3 & \% & 1500 & -- & `drift removed by discarding the largest-magnitude updates` \\
`drift\_spread\_across\_ranks` & 0.079 & -- & 5 & -- & `range of mean drift across assumed ranks 4 to 64` \\
`leakage\_drift\_ratio\_max\_over\_min` & 5.6863 & -- & 200 & -- & `drift at full leakage minus drift at zero, both at norm 0.4` \\
`selection\_efficiency\_keep0\_7` & 0.999 & x & 2000 & -- & `drift removed per unit of update mass removed` \\
`selection\_efficiency\_keep0\_8` & 0.999 & x & 2000 & -- & `drift removed per unit of update mass removed` \\
`selection\_efficiency\_keep0\_9` & 0.999 & x & 2000 & -- & `drift removed per unit of update mass removed` \\
`selection\_efficiency\_keep0\_95` & 0.998 & x & 2000 & -- & `drift removed per unit of update mass removed` \\
`signal\_removed\_keep0\_7` & 65.41 & \% & 2000 & -- & `share of total update magnitude removed by the same rule` \\
`signal\_removed\_keep0\_8` & 52.84 & \% & 2000 & -- & `share of total update magnitude removed by the same rule` \\
`signal\_removed\_keep0\_9` & 35.3 & \% & 2000 & -- & `share of total update magnitude removed by the same rule` \\
`signal\_removed\_keep0\_95` & 22.56 & \% & 2000 & -- & `share of total update magnitude removed by the same rule` \\
\bottomrule
\end{tabular}}
\end{table}

\textbf{40 measurements across 5 artifacts.} Confidence intervals are percentile bootstrap where reported; an em dash marks a quantity that is exact rather than sampled, for which an interval would be meaningless.


\section{Artifact Digests}
\begin{table}[htbp]
\centering
\resizebox{\ifdim\width>\columnwidth\columnwidth\else\width\fi}{!}{%
\begin{tabular}{lc}
\toprule
Artifact & SHA-256 (first 16) \\
\midrule
`artifacts/drift\_by\_norm.json` & `d4d915c4c08b5077` \\
`artifacts/filter\_comparison.json` & `3d81576d3172d485` \\
`artifacts/leakage\_effect.json` & `b55e1e368ef96786` \\
`artifacts/rank\_sensitivity.json` & `266dc78424f564a4` \\
`artifacts/selection\_effect.json` & `b5294b57ed449856` \\
\bottomrule
\end{tabular}}
\end{table}

Any reported value can be recomputed from the artifact named beside it. A digest that no longer matches means the artifact changed after the value was recorded, which invalidates the row rather than the artifact.


\par\vspace{0.5em}
\balance
\bibliographystyle{IEEEtran}
\bibliography{references}

\end{document}
